\documentclass[../DoAn.tex]{subfiles}
\begin{document}

% \section{Thiết kế kiến trúc}
% \subsection{Lựa chọn kiến trúc phần mềm}
% Mục này có độ dài từ một đến ba trang. Sinh viên cần lựa chọn kiến trúc phần mềm cho ứng dụng của mình như: kiến trúc ba lớp MVC, MVP, SOA, Microservice, v.v. rồi giải thích sơ bộ về kiến trúc đó (không giải thích chi tiết/dài dòng).
% Sử dụng kiến trúc phần mềm đã chọn ở trên, sinh viên mô tả kiến trúc cụ thể cho ứng dụng của mình. Gợi ý: sinh viên áp dụng lý thuyết chung vào hệ thống/sản phẩm của mình như thế nào, có thay đổi, bổ sung hoặc cải tiến gì không. Ví dụ, thành phần M trong kiến trúc lý thuyết MVC sẽ là những thành phần cụ thể nào (ví dụ: là interface I + class C1 + class C2, v.v.) trong kiến trúc phần mềm của sinh viên.

% \subsection{Thiết kế tổng quan}
% Sinh viên vẽ biểu đồ gói UML (UML package diagram), nêu rõ sự phụ thuộc giữa các gói (package). SV cần vẽ các gói sao cho chúng được phân theo các tầng rõ ràng, không được sắp đặt package lộn xộn trong hình vẽ. Sinh viên chú ý các quy tắc thiết kế (Các gói không phụ thuộc lẫn nhau, gói tầng dưới không phụ thuộc gói tầng trên, không phụ thuộc bỏ qua tầng, v.v.) và cần giải thích sơ lược về mục đích/nhiệm vụ của từng package. SV tham khảo ví dụ minh họa trong Hình \ref{fig:Fig1}
% \begin{figure}[H]
%     \centering
%     \includegraphics{Hinhve/Picture1.png}
%     \caption{Ví dụ biểu đồ phụ thuộc gói}
%     \label{fig:Fig1}
% \end{figure}
% \subsection{Thiết kế chi tiết gói}
% Sinh viên thiết kế và lần lượt vẽ biểu đồ thiết kế cho từng package, hoặc một nhóm các package liên quan để giải quyết một vấn đề gì đó. Khi vẽ thiết kế gói, sinh viên chỉ cần đưa tên lớp, không cần chỉ ra các thành viên phương thức và thuộc tính. SV tham khảo ví dụ minh họa trong Hình \ref{fig:Fig2}.

% Sinh viên cần vẽ rõ ràng quan hệ giữa các lớp trong biểu đồ. Các quan hệ bao gồm: phụ thuộc (dependency), kết hợp (association), kết tập (aggregation), hợp thành (composition), kế thừa (inheritance), và thực thi (implementation). Các quan hệ này đều đã được minh họa trong \ref{fig:Fig2}.

% Sau khi vẽ hình minh họa, sinh viên cần giải thích ngắn gọn về thiết kế của mình. 

% \begin{figure}[H]
%     \centering
%     \includegraphics{Hinhve/Picture2.png}
%     \caption{Ví dụ thiết kế gói}
%     \label{fig:Fig2}
% \end{figure}

% \section{Thiết kế chi tiết}
% \subsection{Thiết kế giao diện}
% Phần này có độ dài từ hai đến ba trang. Sinh viên đặc tả thông tin về màn hình mà ứng dụng của mình hướng tới, bao gồm độ phân giải màn hình, kích thước màn hình, số lượng màu sắc hỗ trợ, v.v. Tiếp đến, sinh viên đưa ra các thống nhất/chuẩn hóa của mình khi thiết kế giao diện như thiết kế nút, điều khiển, vị trí hiển thị thông điệp phản hồi, phối màu, v.v. Sau cùng sinh viên đưa ra một số hình ảnh minh họa thiết kế giao diện cho các chức năng quan trọng nhất. Lưu ý, sinh viên không nhầm lẫn giao diện thiết kế với giao diện của sản phẩm sau cùng.
% \subsection{Thiết kế lớp}
% Phần này có độ dài từ ba đến bốn trang. Sinh viên trình bày thiết kế chi tiết các thuộc tính và phương thức cho một số lớp chủ đạo/quan trọng nhất của ứng dụng (từ 2-4 lớp). Thiết kế chi tiết cho các lớp khác, nếu muốn trình bày, sinh viên đưa vào phần phụ lục.

% Để minh họa thiết kế lớp, sinh viên thiết kế luồng truyền thông điệp giữa các đối tượng tham gia cho 2 đến 3 use case quan trọng nào đó bằng biểu đồ trình tự (hoặc biểu đồ giao tiếp).
% \subsection{Thiết kế cơ sở dữ liệu}
% Phần này có độ dài từ hai đến bốn trang. Sinh viên thiết kế, vẽ và giải thích biểu đồ thực thể liên kết (E-R diagram). Từ đó, sinh viên thiết kế cơ sở dữ liệu tùy theo hệ quản trị cơ sở dữ liệu mà mình sử dụng (SQL, NoSQL, Firebase, v.v.)

% \section{Xây dựng ứng dụng}
% \subsection{Thư viện và công cụ sử dụng}
% Sinh viên liệt kê các công cụ, ngôn ngữ lập trình, API, thư viện, IDE, công cụ kiểm thử, v.v. mà mình sử dụng để phát triển ứng dụng. Mỗi công cụ phải được chỉ rõ phiên bản sử dụng. SV nên kẻ bảng mô tả tương tự như Bảng \ref{table:my_label}. Nếu có nhiều nội dung trình bày, sinh viên cần xoay ngang bảng.

% \begin{table}[H]
% \centering{}
%     \begin{tabular}{lll}
%         \hline
%         \textbf{Mục đích} & \textbf{Công cụ}       & \textbf{Địa chỉ URL}    \\ \hline
%         IDE lập trình     & Eclipse Oxygen a64 bit & http://www.eclipse.org/ \\ \hline
%         v.v.              & v.v.                   & v.v.                    \\ \hline
%         \end{tabular}
%     \caption{Danh sách thư viện và công cụ sử dụng}
%     \label{fig:my_label}
% \end{table}

% \subsection{Kết quả đạt được}
% Sinh viên trước tiên mô tả kết quả đạt được của mình là gì, ví dụ như các sản phẩm được đóng gói là gì, bao gồm những thành phần nào, ý nghĩa, vai trò?

% Sinh viên cần thống kê các thông tin về ứng dụng của mình như: số dòng code, số lớp, số gói, dung lượng toàn bộ mã nguồn, dung lượng của từng sản phẩm đóng gói, v.v. Tương tự như phần liệt kê về công cụ sử dụng, sinh viên cũng nên dùng bảng để mô tả phần thông tin thống kê này.

% \subsection{Minh họa các chức năng chính}
% Sinh viên lựa chọn và đưa ra màn hình cho các chức năng chính, quan trọng, và thú vị nhất. Mỗi giao diện cần phải có lời giải thích ngắn gọn. Khi giải thích, sinh viên có thể kết hợp với các chú thích ở trong hình ảnh giao diện.

% \section{Kiểm thử}
% Phần này có độ dài từ hai đến ba trang. Sinh viên thiết kế các trường hợp kiểm thử cho hai đến ba chức năng quan trọng nhất. Sinh viên cần chỉ rõ các kỹ thuật kiểm thử đã sử dụng. Chi tiết các trường hợp kiểm thử khác, nếu muốn trình bày, sinh viên đưa vào phần phụ lục.
% Sinh viên sau cùng tổng kết về số lượng các trường hợp kiểm thử và kết quả kiểm thử. Sinh viên cần phân tích lý do nếu kết quả kiểm thử không đạt.
% \section{Triển khai}
% Sinh viên trình bày mô hình và/hoặc cách thức triển khai thử nghiệm/thực tế. Ứng dụng của sinh viên được triển khai trên server/thiết bị gì, cấu hình như thế nào. Kết quả triển khai thử nghiệm nếu có (số lượng người dùng, số lượng truy cập, thời gian phản hồi, phản hồi người dùng, khả năng chịu tải, các thống kê, v.v.)

\section{Thiết kế kiến trúc}
\subsection{Lựa chọn kiến trúc phần mềm}
Hệ thống được xây dựng theo kiến trúc phân tầng kết hợp mô hình ``Simulation--Visualization'' (mô phỏng--hiển thị), trong đó phần mô phỏng giao thông chạy trên SUMO (điều khiển bằng TraCI trong Python) và phần hiển thị/điều khiển chạy trên Unity (C\#). Hai phần giao tiếp với nhau thông qua kết nối TCP socket theo cơ chế truyền thông điệp dạng JSON (có ký hiệu kết thúc gói tin \texttt{<END>} để tách thông điệp trên TCP stream).

Lý do lựa chọn kiến trúc này:
\begin{itemize}
	\item \textbf{Tách biệt trách nhiệm}: SUMO xử lý mô phỏng vi mô (phương tiện, đèn giao thông, người đi bộ), Unity tập trung vào trực quan 3D, giao diện người dùng và tương tác.
	\item \textbf{Dễ mở rộng}: có thể bổ sung thêm loại đối tượng (xe buýt, xe ưu tiên), thêm chỉ số thống kê (thời gian hành trình, mật độ) mà không ảnh hưởng mạnh đến phần còn lại.
	\item \textbf{Tương thích công cụ}: SUMO đã có TraCI hỗ trợ điều khiển/đọc trạng thái theo bước mô phỏng; Unity thuận lợi cho hiển thị thời gian thực.
\end{itemize}

Trong cách ánh xạ theo hướng MVC ở mức khái niệm:
\begin{itemize}
	\item \textbf{Model}: trạng thái mô phỏng trong SUMO (vehicle, traffic light, pedestrian) và các tệp mạng (.net.xml, .rou.xml, .sumocfg).
	\item \textbf{View}: các prefab/manager trong Unity để dựng hình đường, xe, đèn và hiển thị UI.
	\item \textbf{Controller}: lớp giao tiếp TCP và bộ xử lý lệnh (Pause/Resume/End) để đồng bộ mô phỏng.
\end{itemize}

\subsection{Thiết kế tổng quan}
Kiến trúc tổng quan gồm 3 tầng chính:
\begin{enumerate}
	\item \textbf{Tầng mô phỏng (Python + SUMO/TraCI)}: khởi tạo mạng giao thông, chạy mô phỏng theo step, trích xuất trạng thái và ghi thống kê kết quả.
	\item \textbf{Tầng tích hợp (TCP + JSON protocol)}: đóng gói dữ liệu trạng thái (đèn giao thông, xe, người đi bộ), xử lý phân mảnh/ghép gói trên TCP stream bằng marker \texttt{<END>}.
	\item \textbf{Tầng hiển thị (Unity)}: nhận dữ liệu theo thời gian thực, cập nhật/khởi tạo đối tượng trong scene; cung cấp UI cho người dùng điều khiển và kết thúc mô phỏng.
\end{enumerate}

\begin{figure}[htbp]
	\centering
	\fbox{\parbox{0.92\linewidth}{
	\textbf{Hình minh họa kiến trúc tổng quan (tự vẽ theo mô tả):}\\
	Unity (C\#) \textrightarrow[JSON/TCP]{port 5050} Python Main \textrightarrow{TraCI} SUMO.\\
	Unity gửi lệnh điều khiển (Pause/Resume/End) qua port điều khiển, Python phản hồi bằng dữ liệu trạng thái theo bước mô phỏng.
	}}
	\caption{Kiến trúc tổng quan của hệ thống mô phỏng--hiển thị}
	\label{fig:arch_overview}
\end{figure}

\subsection{Thiết kế chi tiết gói}
Các gói/module chính và nhiệm vụ:
\begin{itemize}
	\item \textbf{Python/HelloWorld}:
	\begin{itemize}
		\item \texttt{main.py}: điểm vào chương trình; khởi tạo server socket, nhận lệnh từ Unity, chạy vòng lặp TraCI và gửi dữ liệu trạng thái mỗi step.
		\item \texttt{network.py}: đóng gói giao thức truyền thông; đảm bảo tách thông điệp theo \texttt{<END>} để khắc phục đặc tính stream của TCP.
		\item \texttt{result.py}: ghi log hành trình xe ra CSV và tổng hợp thống kê ra JSON.
		\item \texttt{Traffic/*}: các module đọc trạng thái (xe, đèn, người đi bộ, đường) từ TraCI và chuẩn hóa về dạng JSON.
		\item \texttt{SUMO\_xml/*}: sinh tệp cấu hình/mạng SUMO (node/edge/route/net/sumocfg), trong đó có chức năng sinh mạng từ mê cung.
	\end{itemize}
	\item \textbf{Unity/Assets/Scripts}:
	\begin{itemize}
		\item \texttt{TrafficDataListener}: kết nối tới Python, gửi thông điệp \texttt{SimulationReady} và lắng nghe trạng thái mô phỏng theo thời gian thực.
		\item \texttt{RoadDataListener}: gửi \texttt{RoadDataRequest} để tải dữ liệu đường/nút giao một lần khi khởi động scene.
		\item \texttt{*Manager} (VehicleManager, TrafficLightManager, PedestrianManager): cập nhật/khởi tạo đối tượng dựa trên danh sách dữ liệu nhận được.
		\item \texttt{UI/OptionUI}: UI tùy chọn và gửi lệnh kết thúc mô phỏng (\texttt{Simulation end}).
	\end{itemize}
\end{itemize}

\section{Thiết kế chi tiết}
\subsection{Thiết kế giao diện}
Ứng dụng hướng tới hiển thị mô phỏng thời gian thực trong Unity. Giao diện được thiết kế tối giản, tập trung vào:
\begin{itemize}
	\item \textbf{Màn hình hiển thị chính}: camera quan sát mạng đường, xe và đèn giao thông; các đối tượng được cập nhật theo dữ liệu từ SUMO.
	\item \textbf{Bảng tùy chọn}: bật/tắt UI trong khi chạy, cho phép kết thúc mô phỏng an toàn (gửi lệnh về phía Python).
\end{itemize}
Giả định môi trường chạy phổ biến là PC, độ phân giải khuyến nghị 1920\,$\times$\,1080 (Full HD) hoặc cao hơn; UI được thiết kế theo nguyên tắc dễ đọc, hạn chế chữ nhỏ và đảm bảo thao tác bằng chuột/bàn phím.

\subsection{Thiết kế lớp}
Trong phạm vi báo cáo, trình bày chi tiết một số lớp/module chủ đạo.

\subsubsection*{Nhóm lớp Python (mô phỏng và ghi log)}
\textbf{VehicleTripCsvLogger} (trong \texttt{result.py}):
\begin{itemize}
	\item Trách nhiệm: ghi nhận thời điểm xe xuất phát/đến đích theo step mô phỏng; ghi CSV và xuất thống kê trung bình ra JSON.
	\item Thuộc tính chính: \texttt{\_depart\_step\_by\_vehicle}, \texttt{\_travel\_steps\_sum}, \texttt{\_travel\_steps\_count}.
	\item Phương thức chính: \texttt{open()}, \texttt{on\_step(step, departed, arrived)}, \texttt{close()}.
\end{itemize}

\textbf{network.py} (giao thức TCP):
\begin{itemize}
	\item Trách nhiệm: tách/gộp thông điệp trên TCP bằng marker \texttt{<END>}; tránh lỗi khi một lần \texttt{recv} chứa nhiều thông điệp hoặc chỉ chứa một phần thông điệp.
	\item Hàm chính: \texttt{send\_data()}, \texttt{receive\_data()}, \texttt{receive\_message()}, \texttt{server\_thread()}.
\end{itemize}

\subsubsection*{Nhóm lớp Unity (nhận dữ liệu và cập nhật scene)}
\textbf{TrafficDataListener}:
\begin{itemize}
	\item Trách nhiệm: kết nối TCP đến Python, gửi \texttt{SimulationReady}, đọc JSON trạng thái; đẩy xử lý sang main thread để cập nhật scene an toàn.
	\item Phương thức chính: \texttt{StartListening()}, \texttt{ListenForData()}, \texttt{HandleData(json)}.
\end{itemize}

\textbf{RoadDataListener}:
\begin{itemize}
	\item Trách nhiệm: gửi yêu cầu \texttt{RoadDataRequest} và nhận dữ liệu đường/nút giao 1 lần; lưu vào \texttt{RoadDataSO} để các hệ thống khác sử dụng.
\end{itemize}

\textbf{TrafficLightManager}:
\begin{itemize}
	\item Trách nhiệm: quản lý tập đối tượng đèn giao thông trong scene; nếu đã tồn tại thì cập nhật trạng thái/vị trí, nếu chưa thì tạo mới từ prefab.
\end{itemize}

\subsubsection*{Luồng thông điệp (mô tả dạng trình tự)}
\textbf{Use case 1 --- Unity yêu cầu dữ liệu đường (RoadData)}:
\begin{enumerate}
	\item Unity khởi chạy \texttt{RoadDataListener} và mở kết nối TCP tới Python.
	\item Unity gửi chuỗi \texttt{RoadDataRequest}.
	\item Python trả về JSON mô tả edge/junction/crossing (đóng gói và kết thúc bằng \texttt{<END>}).
	\item Unity parse JSON và lưu vào \texttt{RoadDataSO}; sau đó dừng listener.
\end{enumerate}

\textbf{Use case 2 --- Đồng bộ trạng thái mô phỏng theo thời gian thực}:
\begin{enumerate}
	\item Unity gửi \texttt{SimulationReady} để bắt đầu vòng lặp mô phỏng.
	\item Python chạy TraCI theo step; mỗi step đọc danh sách xe/đèn/người đi bộ và gửi về Unity.
	\item Unity nhận JSON, phân rã thành danh sách và gọi các manager để cập nhật đối tượng trong scene.
	\item Khi người dùng thoát ứng dụng, Unity gửi \texttt{Simulation end} để yêu cầu dừng mô phỏng.
\end{enumerate}

\subsection{Thiết kế cơ sở dữ liệu}
Hệ thống không sử dụng cơ sở dữ liệu theo nghĩa truyền thống. Dữ liệu được lưu dưới dạng tệp:
\begin{itemize}
	\item \textbf{Dữ liệu cấu hình/mạng mô phỏng}: các tệp XML của SUMO (\texttt{.nod.xml}, \texttt{.edg.xml}, \texttt{.net.xml}, \texttt{.rou.xml}, \texttt{.sumocfg}).
	\item \textbf{Dữ liệu kết quả thực nghiệm}: CSV ghi thời gian hành trình và JSON tổng hợp thống kê theo lần chạy (trong thư mục \texttt{result/}).
\end{itemize}
Thiết kế này phù hợp vì mục tiêu là mô phỏng/hiển thị và thu thập thống kê theo từng lần chạy, không cần truy vấn phức tạp hay đa người dùng.

\section{Xây dựng ứng dụng}
\subsection{Thư viện và công cụ sử dụng}
Bảng \ref{table:tools} liệt kê các công cụ/chính sách phiên bản sử dụng trong dự án.

\begin{table}[htbp]
\centering
\begin{tabular}{p{0.35\linewidth}p{0.28\linewidth}p{0.28\linewidth}}
\hline
\textbf{Mục đích} & \textbf{Công cụ} & \textbf{Phiên bản/Ghi chú} \\ \hline
Mô phỏng giao thông vi mô & Eclipse SUMO & 1.24.0 (tạo file \texttt{.sumocfg} thể hiện version) \\ \hline
Điều khiển SUMO & TraCI (Python API) & đi kèm SUMO 1.24.0 \\ \hline
Ngôn ngữ mô phỏng/tích hợp & Python & (điền theo máy: chạy \texttt{python --version}) \\ \hline
Trực quan 3D + UI & Unity Editor & 6000.3.0f1 \\ \hline
Parse JSON (Unity) & Newtonsoft.Json & com.unity.nuget.newtonsoft-json 3.2.2 \\ \hline
Môi trường phát triển & Visual Studio Code & (điền theo máy: About VS Code) \\ \hline
Quản lý phiên bản & Git & (điền theo máy: \texttt{git --version}) \\ \hline
Hệ điều hành chạy thử nghiệm & Windows & (10/11 tùy máy triển khai) \\ \hline
\end{tabular}
\caption{Danh sách thư viện và công cụ sử dụng}
\label{table:tools}
\end{table}

\subsection{Kết quả đạt được}
Các sản phẩm/dữ liệu đầu ra chính:
\begin{itemize}
	\item \textbf{Bộ sinh bản đồ/mạng SUMO}: tạo node--edge từ bản đồ thành phố hoặc từ mê cung; sinh ra các tệp \texttt{.nod.xml}, \texttt{.edg.xml}, \texttt{.net.xml}, \texttt{.rou.xml}, \texttt{.sumocfg}.
	\item \textbf{Ứng dụng Unity hiển thị mô phỏng}: nhận dữ liệu thời gian thực từ Python và cập nhật đối tượng (xe, đèn, người đi bộ).
	\item \textbf{Bộ ghi log kết quả}: xuất CSV hành trình xe theo step và JSON tổng hợp theo từng lần chạy.
\end{itemize}

Thống kê nhanh mã nguồn (tính trên các phần do nhóm tự viết, loại trừ thư viện/\texttt{.venv}):
\begin{itemize}
	\item Python: 20 file, 3045 dòng, dung lượng khoảng 109.5\,KB.
	\item Unity C\# (Assets/Scripts): 72 file, 4261 dòng, dung lượng khoảng 126.9\,KB.
	\item Tệp cấu hình/mạng SUMO (HelloWorld/SUMO\_xml): 6 tệp, tổng dung lượng khoảng 10.34\,MB.
	\item Bản build Windows minh họa: \texttt{UnityBuild/TestGR1.1.exe} khoảng 0.64\,MB (chưa tính toàn bộ thư mục build).
\end{itemize}

\subsection{Minh họa các chức năng chính}
Các chức năng chính đã thực hiện:
\begin{enumerate}
	\item \textbf{Sinh mạng đường từ dữ liệu đầu vào}: hỗ trợ sinh mạng từ tệp bản đồ thành phố và sinh mạng từ mê cung (tạo node/edge theo cấu trúc lưới).
	\item \textbf{Chạy mô phỏng SUMO bằng TraCI}: điều khiển mô phỏng theo step và đọc trạng thái xe/đèn/người đi bộ.
	\item \textbf{Đồng bộ dữ liệu sang Unity theo thời gian thực}: gửi gói JSON theo TCP; Unity nhận và cập nhật scene.
	\item \textbf{Điều khiển mô phỏng}: hỗ trợ tạm dừng/tiếp tục và kết thúc mô phỏng từ Unity.
	\item \textbf{Xuất kết quả thực nghiệm}: ghi log hành trình xe và thống kê trung bình theo lần chạy.
\end{enumerate}

\begin{figure}[htbp]
	\centering
	\fbox{\parbox{0.92\linewidth}{\textbf{Chèn ảnh minh họa 1:} Scene Unity đang hiển thị xe/đèn giao thông theo dữ liệu SUMO.}}
	\caption{Minh họa trực quan mô phỏng trên Unity}
	\label{fig:demo_unity}
\end{figure}

\begin{figure}[htbp]
	\centering
	\fbox{\parbox{0.92\linewidth}{\textbf{Chèn ảnh minh họa 2:} SUMO-GUI đang chạy mô phỏng và cấu hình delay (ví dụ 110ms) để đồng bộ quan sát.}}
	\caption{Minh họa mô phỏng trên SUMO-GUI}
	\label{fig:demo_sumo}
\end{figure}

\section{Kiểm thử}
Do dự án có tính tích hợp cao (SUMO \& Unity), kiểm thử được thực hiện theo hai hướng: (i) kiểm thử chức năng thủ công theo use case và (ii) kiểm thử hộp đen cho các module xử lý dữ liệu/giao thức.

\subsection*{Các trường hợp kiểm thử tiêu biểu}
\begin{table}[htbp]
\centering
\begin{tabular}{p{0.12\linewidth}p{0.33\linewidth}p{0.45\linewidth}}
\hline
\textbf{TC} & \textbf{Mục tiêu} & \textbf{Dữ liệu vào / Kỳ vọng} \\ \hline
TC01 & Tách thông điệp TCP đúng theo \texttt{<END>} & Gửi liên tiếp 2 JSON trong 1 lần \texttt{sendall}; kỳ vọng phía nhận tách đúng 2 thông điệp, không lẫn chuỗi. \\ \hline
TC02 & Nhận RoadData một lần & Unity gửi \texttt{RoadDataRequest}; kỳ vọng nhận JSON hợp lệ và lưu vào \texttt{RoadDataSO}, sau đó listener dừng. \\ \hline
TC03 & Cập nhật trạng thái đèn giao thông & SUMO trả về chuỗi trạng thái \texttt{r/y/g/G}; kỳ vọng Unity ánh xạ về trạng thái hiển thị tương ứng (đỏ/vàng/xanh). \\ \hline
\end{tabular}
\caption{Một số trường hợp kiểm thử chính}
\label{table:testcases}
\end{table}

Kết quả: các ca kiểm thử trên đạt yêu cầu trong điều kiện chạy thử nghiệm. Trường hợp sai lệch đồng bộ (giật/nhảy vị trí) chủ yếu xuất hiện khi delay quan sát trong SUMO-GUI đặt quá thấp; khi tăng delay (ví dụ khoảng 110ms) thì quan sát ổn định hơn.

\section{Triển khai}
\subsection*{Môi trường triển khai}
Hệ thống được triển khai thử nghiệm trên máy cá nhân Windows với 2 tiến trình chính:
\begin{itemize}
	\item \textbf{Python + SUMO}: chạy script chính để khởi tạo server và điều khiển TraCI.
	\item \textbf{Unity (Editor hoặc bản build)}: chạy scene nhận dữ liệu và hiển thị.
\end{itemize}

\subsection*{Quy trình chạy thử nghiệm}
Quy trình triển khai thử nghiệm (tóm tắt theo README):
\begin{enumerate}
	\item Thêm dự án Unity vào Unity Hub và mở project.
	\item Chạy script Python \texttt{main.py} để chuẩn bị kết nối và khởi chạy mô phỏng.
	\item Khi SUMO-GUI mở, đặt delay quan sát phù hợp (ví dụ 110) và chạy mô phỏng.
	\item Chạy Unity (Play mode hoặc bản build) để bắt đầu nhận dữ liệu và hiển thị.
\end{enumerate}

Các cổng TCP sử dụng theo cấu hình hiện tại:
\begin{itemize}
	\item Port 5050: Unity kết nối để nhận dữ liệu RoadData/TrafficData.
	\item Port điều khiển (ví dụ 5054): Unity gửi lệnh kết thúc mô phỏng (\texttt{Simulation end}).
\end{itemize}

Trong phạm vi thử nghiệm, hệ thống chạy ổn định theo mô hình ``1 máy -- 1 người dùng''; dữ liệu kết quả được lưu cục bộ thành CSV/JSON để phục vụ phân tích.

\end{document}
